\documentclass[../../../include/open-logic-section]{subfiles}

\begin{document}

\olfileid{sth}{choice}{hartogs}
\olsection{તુલનીયતા અને હાર્ટોગ્સનું લેમા}

આ તેની સારી બાજુ છે. હવે ખરાબ બાજુ જુઓ.
પસંદગી વિના બાબતો \emph{જટિલ} બને છે.
એ શા માટે, તે સમજવા \cite{Hartogs1915}નું
એક સુંદર પરિણામ જોઈએ:

\begin{lem}[\emph{$\ZF$માં}]\ollabel{HartogsLemma}
કોઈ પણ ગણ $A$ માટે એવી ક્રમવાચક સંખ્યા
$\alpha$ છે કે $\cardnless{\alpha}{A}$.
\end{lem}

\begin{proof}
જો $B \subseteq A$ અને $R \subseteq B^2$ હોય,
તો \olref[ord-arithmetic][using-addition]{rankcomputation}
મુજબ $\tuple{B, R} \subseteq
V_{\setrank{A}+4}$. તેથી પૃથક્કરણથી આ ગણ લો:
\[
	C = \Setabs{\tuple{B, R} \in V_{\setrank{A}+5}}{B\subseteq A 
	\text{ અને $\tuple{B, R}$ સુક્રમ છે}}
\]
પ્રતિસ્થાપન અને
\olref[ordinals][ordtype]{thmOrdinalRepresentation}
વાપરી આ ગણ રચો:
\[
	\alpha = \Setabs{\ordtype{B, R}}{\tuple{B, R} \in C}.
\]
\olref[ordinals][basic]{corordtransitiveord} મુજબ
$\alpha$ ક્રમવાચક સંખ્યા છે, કારણ કે તે
ક્રમવાચક સંખ્યાઓનો પરંપરિત ગણ છે.
ખરેખર $\gamma \in \beta \in \alpha$ હોય,
તો કોઈ $B \subseteq A$ માટે
$\beta = \ordtype{B, R}$ છે. પછી
\olref[ordinals][iso]{wellordinitialsegment}
મુજબ કોઈ $b \in B$ માટે
$\gamma = \ordtype{B_b, R_b}$, તેથી
$\gamma \in \alpha$.
\sourcecorrection{OLMD-144}{મૂળ પુરાવામાં અહીં
$B \subseteq R$ લખ્યું છે, જ્યારે $R$ એ
$B$ પરનો સંબંધ છે અને રચનામાં જરૂરી
ધારણા $B \subseteq A$ છે.}

વિરોધાભાસ માટે ધારો કે !!a{injection}
$f \colon \alpha \to A$ છે. હવે આ રીતે લો:
\begin{align*}
	B &= \ran{f}\\
	R &= \Setabs{\tuple{f(\gamma), f(\delta)} \in A \times A}{\gamma,\delta\in\alpha \land \gamma \in \delta}.
\end{align*}
સ્પષ્ટપણે $\alpha = \ordtype{B, R}$ અને
$\tuple{B, R} \in C$. તેથી $\alpha
\in \alpha$, જે વિરોધાભાસ છે.
\sourcecorrection{OLMD-145}{મૂળ સંબંધ-સૂત્રમાં
$f(\alpha)$ છે, પરંતુ $f$નું ક્ષેત્ર $\alpha$
હોવાથી $f(\alpha)$ વ્યાખ્યાયિત નથી.
સંબંધની જોડી માટે $\gamma,\delta\in\alpha$
એવા નવા સૂચકો લીધા છે અને
$\gamma\in\delta$ શરત સ્પષ્ટ કરી છે.}
\end{proof}

આમાંથી એક ઊંડું પરિણામ મળે છે:

\begin{thm}[\emph{$\ZF$માં}]
નીચેના દાવા સમકક્ષ છે:
\begin{enumerate}
	\item\ollabel{equivwo} સુક્રમિતતાની સ્વયંસિદ્ધિ
	\item\ollabel{equivcompare} કોઈ પણ ગણો $A$ અને $B$ માટે
	$\cardle{A}{B}$ અથવા $\cardle{B}{A}$
\end{enumerate}
\end{thm}

\begin{proof}
\emph{\olref{equivwo} $\Rightarrow$ \olref{equivcompare}.}
$A$ અને $B$ નિશ્ચિત લો. \olref{equivwo}થી
$\tuple{A, R}$ અને $\tuple{B, S}$ સુક્રમો છે.
\olref[ordinals][ordtype]{thmOrdinalRepresentation}
વાપરી $f \colon \alpha \to \tuple{A, R}$
અને $g \colon \beta \to \tuple{B, S}$
સમરૂપતાઓ લો.
\olref[sth][ordinals][basic]{ordinalsaresubsets}
મુજબ $\alpha \subseteq \beta$ અથવા
$\beta \subseteq \alpha$. જો
$\alpha \subseteq \beta$, તો
$\comp{f^{-1}}{g} \colon A \to B$
એ !!a{injection} છે, એટલે
$\cardle{A}{B}$; તેવી જ રીતે
$\beta \subseteq \alpha$ હોય તો
$\cardle{B}{A}$.

\emph{\olref{equivcompare} $\Rightarrow$ \olref{equivwo}.}
$A$ નિશ્ચિત લો. \olref{HartogsLemma}થી
એવી ક્રમવાચક સંખ્યા $\beta$ છે કે
$\cardnless{\beta}{A}$. હવે
\olref{equivcompare}થી
$\cardle{A}{\beta}$. તેથી !!{injection}
$f \colon A \to \beta$ છે, અને તેના વડે
$A$ના ઘટકોને સુક્રમિત કરવા આ ક્રમ
વ્યાખ્યાયિત કરી શકીએ:
$\Setabs{\tuple{a, b} \in A \times A}{f(a)
\in f(b)}$.
\end{proof}
\noindent
તરત જ એ પરિણામ મળે છે કે સુક્રમિતતા
નિષ્ફળ જાય, તો કેટલાક ગણો કદની દૃષ્ટિએ
\emph{ખરેખર અતુલનીય} હોય છે. એટલે
સુક્રમિતતા નિષ્ફળ જાય, તો અનંતાતીત
ગણાંક-અંકગણિત જટિલ બનશે. ઉદાહરણ તરીકે,
$A$ અને $B$ અનંત હોય, તો
$\cardeq{A \disjointsum B}{M}$ અને
$\cardeq{A \times B}{M}$, જ્યાં
$M$ એ $A$ અને $B$માંથી મોટો ગણ છે,
એવો વિચાર છોડવો પડશે
(\olref[card-arithmetic][simp]{cardplustimesmax} જુઓ).
સમસ્યા સરળ છે: જો $A$ અને $B$ના કદની
\emph{તુલના} જ કરી ન શકીએ, તો કયો મોટો
છે તે પૂછવાનો અર્થ નથી.
\sourcecorrection{OLMD-146}{મૂળના અંતિમ સૂત્રમાં
$\cardeq$ને ફરી $\cardeq$ના પહેલા સ્થાનમાં
મૂક્યું છે, એટલે તે ગણોની સમકદતાનો
દાવો વ્યક્ત કરતું નથી. અહીં ઇચ્છિત બે
સમકદતા-દાવા અલગ લખ્યા છે.}

\end{document}
